
\documentclass{article}
\usepackage{colortbl}
\usepackage{makecell}
\usepackage{multirow}
\usepackage{supertabular}

\begin{document}

\newcounter{utterance}

\centering \large Interaction Transcript for game `wordle\_withclue', experiment `high\_frequency\_words\_clue\_no\_critic', episode 1 with reflexion{-}Qwen3.5{-}9B.
\vspace{24pt}

{ \footnotesize  \setcounter{utterance}{1}
\setlength{\tabcolsep}{0pt}
\begin{supertabular}{c@{$\;$}|p{.15\linewidth}@{}p{.15\linewidth}p{.15\linewidth}p{.15\linewidth}p{.15\linewidth}p{.15\linewidth}}
    \# & \multicolumn{2}{c}{Player} && \multicolumn{2}{c}{Game Master} \\
    \hline

    \theutterance \stepcounter{utterance}  
    & & & \multicolumn{4}{p{0.6\linewidth}}{
        \cellcolor[rgb]{0.9,0.9,0.9}{
            \makecell[{{p{\linewidth}}}]{
                \texttt{\tiny{[P1$\langle$GM]}}
                \texttt{You have attempted this same task before and did not succeed. These are the plans you wrote after those attempts.} \\
\texttt{Trial \#1: New plan: I failed because I guessed the clue word "comatose" directly, which is likely incorrect as the clue is meant to hint at the target, not be the target itself. Additionally, my explanation was too verbose and did not strictly adhere to the required format of providing a concise explanation followed immediately by the guess. Next time, I will ignore the literal meaning of the clue as the answer, instead using it to generate a plausible 5{-}letter word that fits the theme (e.g., if the clue is "comatose," I might guess a related state like "sleep" or a synonym like "drows" if valid, but strictly sticking to 5 letters). I will ensure my response contains only the "explanation:" tag with a brief reasoning linking the clue to the chosen word, followed by the "guess:" tag with a valid 5{-}letter word, avoiding any extra text or premature conclusions.} \\
\\ 
\texttt{Trial \#2: New plan: I failed because my explanation became a rambling internal monologue that violated the "Do not generate any other text" rule, causing the Game Master to abort the turn. Additionally, I failed to synthesize the feedback from "sleep" (e at 4) and "dazed" (d at 1 or 5, e at 4) into a concrete word pattern before guessing. Next time, I will strictly limit my explanation to a single sentence summarizing the letter constraints derived from previous feedback and the clue, immediately followed by a valid 5{-}letter word that fits the pattern (e.g., if 'd' is yellow and 'e' is green at 4, I will deduce a word like "dread" or "dosed" if they fit, or pivot to a new pattern if none do), ensuring no extra text or incomplete thoughts are included.} \\
\\ 
\texttt{Trial \#3: New plan: I failed because I entered a repetitive loop by guessing "dazed" multiple times after realizing 'd' was not at the ends and 'e' was fixed at position 4, without exploring other valid 5{-}letter words that fit the pattern `\_ \_ \_ e \_` and the "comatose" theme (e.g., "stale", "fused", "mused", "haste"). My strategy lacked a systematic method to generate new candidates once a letter was ruled out from specific positions. Next time, upon receiving feedback that 'e' is at position 4 and 'd' is elsewhere, I will immediately list all valid 5{-}letter words ending in "e" (or fitting `\_ \_ \_ e \_`) that match the clue's semantic field, test a new candidate from this list (like "stale" or "fused") instead of reusing the previous guess, and ensure my explanation explicitly states the new word choice is derived from this expanded search rather than just restating the feedback.} \\
\\ 
\\ 
\texttt{You are a language wizard who likes to guess words by using the given rules.} \\
\\ 
\texttt{Welcome to Wordle! You have six attempts to guess the target word, a valid English word of five lowercase letters (a{-}z). Please use the tags "explanation:" and "guess:" to provide a concise explanation for each guess.} \\
\\ 
\texttt{To help you make an informed guess, you will receive a clue for the word, such as} \\
\texttt{clue: snowy white.} \\
\\ 
\texttt{Here is an example guess based on the clue:} \\
\texttt{explanation: In the fairy tail Snow White, the girl is killed because she eats a poisoned apple. And the word apple has 5 letters.} \\
\texttt{guess: apple} \\
\\ 
\texttt{After each guess, your answer will be validated, and you will receive feedback indicating which letters are correct (green), which letters are correct but in the wrong position (yellow), and which letters are incorrect (red). This feedback can be useful in determining which letters to include or exclude in your next guess.} \\
\\ 
\texttt{For example, the feedback for "apple" might be:} \\
\texttt{guess\_feedback: a<yellow> p<yellow> p<green> l<yellow> e<red>} \\
\\ 
\texttt{The explanation should contain details about how both the guess\_feedback and clue work together in tandem to formulate a new guess.} \\
\\ 
\texttt{Do not generate any other text.} \\
\\ 
\texttt{Let's begin with your first guess.} \\
\\ 
\texttt{clue: comatose} \\
            }
        }
    }
     \\ \\

    \theutterance \stepcounter{utterance}  
    & \multicolumn{4}{p{0.6\linewidth}}{
        \cellcolor[rgb]{0.561,0.737,0.561}{
            \makecell[{{p{\linewidth}}}]{
                \texttt{\tiny{[P1$\rangle$GM]}}
                \texttt{explanation: The clue "comatose" suggests a state of deep sleep or unconsciousness; I will guess a 5{-}letter word related to this state, such as "slumber" or "drowse", but since I need a valid 5{-}letter word that fits the theme, I'll try "drowse" as it directly relates to being comatose.} \\
\texttt{guess: drowse} \\
            }
        }
    }
    & & \\ \\

    \theutterance \stepcounter{utterance}  
    & & & \multicolumn{2}{p{0.3\linewidth}}{
        \cellcolor[rgb]{0.9,0.9,0.9}{
            \makecell[{{p{\linewidth}}}]{
                \texttt{\tiny{[GM$|$GM]}}
                \texttt{Error: The length of the guessed word is not 5.} \\
            }
        }
    }
    & & \\ \\

    \theutterance \stepcounter{utterance}  
    & & & \multicolumn{2}{p{0.3\linewidth}}{
        \cellcolor[rgb]{0.9,0.9,0.9}{
            \makecell[{{p{\linewidth}}}]{
                \texttt{\tiny{[GM$|$GM]}}
                \texttt{game\_result = ABORT} \\
            }
        }
    }
    & & \\ \\

    \theutterance \stepcounter{utterance}  
    & & \multicolumn{4}{p{0.6\linewidth}}{
        \cellcolor[rgb]{0.95,0.95,0.95}{
            \makecell[{{p{\linewidth}}}]{
                \texttt{\tiny{[P1$\leftrightarrow$P]}}
                \texttt{Post{-}episode Reflexion (not sent to GM):} \\
\texttt{New plan: I failed because my explanation contained internal monologue and conditional reasoning ("such as...", "but since...") which violated the strict "Do not generate any other text" rule and the requirement for a concise, direct explanation. Additionally, I did not explicitly link the clue "comatose" to the specific word chosen in the explanation. Next time, I will strictly format my response as two distinct lines: first, a single sentence under "explanation:" that directly states the clue's meaning and identifies the chosen 5{-}letter word as a synonym or related concept (e.g., "The clue 'comatose' indicates deep sleep, so I will guess 'slumber'"), and immediately follow it with "guess:" and the valid 5{-}letter word, ensuring no extra text, brackets, or hypothetical scenarios are included.} \\
            }
        }
    }
    & \\ \\

\end{supertabular}
}

\end{document}
